The update-latency measurements reported in Experiment~5 reflect single-process behavior: one U-HAP instance compiles the updated policy set into indexed artifacts and activates them via an atomic registry pointer-swap, replacing the live index in a single compare-and-exchange instruction with no request-processing interruption.
Production Kubernetes deployments run multiple webhook replicas behind a load balancer, so the practical consistency window is not a single swap but the maximum per-replica swap latency across all replicas~$\max_{i}\,\delta_i$, where each~$\delta_i$ is the compilation-plus-swap time measured here.
Because U-HAP's Phase~2 compilation is CPU-bound and deterministic, all replicas process the same policy bundle independently and converge to the identical compiled state; the consistency window closes once the slowest replica completes its swap, typically within the same order of magnitude as the single-process figure.
Per-replica swap latency is therefore the correct primitive for characterizing update overhead: it directly bounds the consistency window and determines how quickly any individual replica stops serving stale decisions.
Horizontal-coordination mechanisms such as leader-election-gated bundle activation or a two-phase distributed commit are orthogonal to U-HAP's compilation-driven design and can be layered atop the existing registry interface without altering Phase~2 or Phase~3 logic; we leave their empirical evaluation to future work on large-scale multi-tenant deployments.
